\section{The Hodge words at the inert primes}

% BEGIN I82W1
\paragraph{I82W1. A single cutoff for inert primes and their size parameters.}
Use the original family and the hub of F78D2. Fix its row constants
$R$ and write $C_2=R.C_2$, $C=\max(100,2C_2)$, $W=w''(n)$,
$S=S_c$, $\eta=\eta_{\rm Sz}$. All assertions below hold after
one cutoff in the family variable $x$, before the choice of the
member, index, point, completion or coefficient field. Indeed the
positive exponents $\epsilon,\epsilon_z$ are family constants (use
a member at zero); for $x>x_1$ the parameter tie fixes the member's
$C_{L_w}$ at the eighth entry of the family list, with default zero.
The parameter inequalities give $t\ge7$, $t\le\eta$,
$x t\eta\le W$, $kS\le W$, $k\ge m\sqrt{x}\log x$ with $m\ge2$,
and $W=b'/(xn)\le b'$. At $x\ge10^4$ they imply
\[
 W\ge100,\qquad 100(\eta+S+t+2)\le W .
\]
For a top place the upper center bound is
$b_\tau\le b'(1+T)$ with $T=N_\star x^\epsilon\Lambda_G/k$.
The parameter estimate for $\rho_P$ includes $4T\le x^{-2}$
(the other summand is nonnegative); thus $b_\tau+b'/x\le3b'/2$.
Finally $L_w\le C_{L_w}\log x$ and
$b'/b=x\exp(-x^{\epsilon/3})$. Hence
$C(L_w+1)b'\le b$ eventually, uniformly in all members.
These inequalities are precisely the six clauses of $\mathrm{Roomy}(L,C)$.

At an inert index $p=\ell_i$, the host classification gives
$p\notin{\rm small}$, $p\ne q,d$, and $(-d/p)\ne1$.
The host splitting condition makes the symbol $+1$ for every
odd prime up to its split bound, which is at least
$\lfloor x^{\epsilon_z/5}\rfloor$. Thus $p>x^{\epsilon_z/5}$,
and a uniform increase of the cutoff gives $p\ge11$.
The two type primes are $1$ modulo $4d$ and have symbol $+1$,
so this $p$ is away from both. As in M81B1, $p$ is prime to
the host conductor and every completion $K=E_w$ above it is
unramified: the conductor cyclotomic extension is unramified here.
Write $n_w=[K:\mathbb Q_p]=[K_0:\mathbb Q_p]$.
The degree on the bridge is also the residue degree of $w$,
by the completion degree formula $[E_w:\mathbb Q_p]=e_w f_w$.
% END I82W1

% BEGIN I82W2
\paragraph{I82W2. The letters of the actual representation and their orientation.}
Fix such an $i$, $\pi\in Z_{\rm good}(i)$, $w\mid p$, and any
coefficient extension $L'$ good at level one on this bridge.
DFD7--DFD9 give $\mathrm{GoodR}(X.r_i)$ and the continuous
complement $X.\rho_i(\pi)$, with its fraction-field splitting
at $\pi\in Z(i)$. S80J3 gives crystallinity for every point of
the localized family after extension of coefficients. These are the
goodness, continuity, membership, splitting and crystalline clauses of
$\mathrm{PtIn}$ at the indicated family and point. S80J1 supplies
one orientation $\varepsilon\in\{1,-1\}$ for the member, with
its actual labelled jump readings.

Goodness splits $K$ over $L'$ and includes the conjugates of
$K_0(\mu_p)$. The ring of integers of $L'$ is a DVR with fraction
field $L'$. Since $K=K_0$, for each Wach letter
$a\in\mathbb Z/n_w\mathbb Z$ there is precisely one embedding
$\tau_a:K\to L'$ restricting to that letter. For clarity on the
labelling, write $c:E\to K$ for the dense completion map and choose
a base embedding $\kappa:E\to L'$ so that
$\tau_0 c=\kappa\sigma_{g_0}$ for the host label read at zero.
One may take $\kappa=\tau_0 c\,\sigma_{g_0}^{-1}$: the prescribed
complex bridge reading of $\kappa$ is then the base reading, so all
other labels are read with this same $\kappa$. The next Wach
embedding satisfies $\tau_{a+1}\sigma=\tau_a$, where $\sigma$
is arithmetic Frobenius on $K_0$. Its restriction to the host is
$f_p$, the specified arithmetic Frobenius. Consequently
\[
 \tau_a c=\kappa\sigma_{g_a},\qquad
 g_a=g_0\psi^{\,a},\qquad \psi=f_p^{-1}.
\]
The decomposition group of the unramified completion is cyclic
generated by $f_p$ and has order $n_w$; this is the exact length
of the orbit. Thus the choice of $\kappa$ is independent of $a$.

Put $A=(t+1)/2$, an integer, and let $E_a=e_{j_i-1}$ at the
real-place label of $g_a$ (the common place of $g_a$ and
$\mathrm{cc}\,g_a$). If the CM sign of $g_a$ equals $\varepsilon$,
declare the letter down with size $s_a=E_a+A$; otherwise declare
it up with size $s_a=E_a-A$. Set
\[
 B=\lceil b\rceil-A,\quad
 H=\lfloor\eta\rfloor+t+2,\quad S_0=\lfloor S\rfloor,\quad
 W_0=\max_{\text{nonthick }g,\ j}\lfloor e_j(g)\rfloor+A,
\]
with zero for the maximum over an empty set. The $e_j$ are integral.
In S80J1 the original upward cloud is $(t-1)/2+E_a+v$.
Its twisted dual sends $z$ to $-1-z$ and has cloud
$-s_a-v$. Reflecting the original cloud before dualizing gives
$s_a+v$ at the up letter instead. Here $v$ is integral in $[0,S_0]$.
The lows and the dual string lie in $[-H,H]$. Every lattice
$T_1$ tested in $\mathrm{WordBase}$ embeds rationally into a
finite direct sum of the localized full twisted duals. Crystalline
representations are closed under subobjects in the tied comparison;
the induced labelled de Rham filtration is strict, so graded
multiplicities of a subobject are bounded by those of the sum.
Thus at $a$ the jumps of $T_1$ lie in $[-H,H]$ together with
$[s_a,s_a+S_0]$ at an up letter or $[-s_a-S_0,-s_a]$ at a down
letter. The endpoints of the containing interval $[-M_a,h_a]$
are $(M_a,h_a)=(H,s_a+S_0)$ up and $(s_a+S_0,H)$ down,
using $s_a>H$ proved below. Every tested $T_2$ is rationally
the twisted dual of the rank-$r$ complement at $\pi$.
It is crystalline and its jump multiset has the same cloud and
only the $m$ lows, with no character string. The cloud lies
outside $[-H,H]$; summing the multiplicities there gives at most
$m$. Each letter has just the embedding $\tau_a$, so this is both
the single-embedding multiplicity bound and the full
$\mathrm{lowCount}$ bound.
% END I82W2

% BEGIN I82W3
\paragraph{I82W3. Bounds for both kinds of letter.}
The thick host window gives $b\le E_a\le b+\eta+1$.
Since $E_a$ is integral,
$\lceil b\rceil\le E_a\le\lceil b\rceil+\lfloor\eta\rfloor+1$.
Hence $B\le s_a\le B+H$ at a thick place, including either sign
and the endpoint allowance of one. At a nonthick place the thin
or top window gives
\[
 W\le E_a\le w_B=
 \max\{W+\eta,\ \sup_\tau b_\tau+b'/x+\eta\},
 \qquad E_a\le 8b'/5.
\]
Indeed the thin interval starts at $W$; the lower top center is
at least $W$, the upper top center is at most $3b'/2$, and
$\eta\le b'/100$ by I82W1. Therefore $W_0\le w_B+A=w_P$,
and $W_0+S_0+H+1\le2b'$ by the $8b'/5$ bound and I82W1.
Also $\lvert s_a-E_a\rvert=A$ and $s_a\le W_0$ here.
Since $b\ge100b'$, $A\le b'/100$ and $B\ge b-A$,
\[
 20W_0,\ 20S_0,\ 20H\ \le B,\qquad H<s_a
\]
for every letter (use $s_a\ge W-A$). A thick letter of the word
is consequently exactly one over a thick host place.
The common bounds now give $S_0\le S$, $H\le\eta_3$,
$B+H+S_0\le h_B=b+S+\eta+t+2$, and the word's $m$ is
the original $m$. In fact $S_0\ge1$: the cloud has at least two
distinct nonnegative integral entries $v$ in $[0,S]$, since
$k\ge m+3$. Also $H,W_0,B\ge1$. In the unfrozen
$\mathrm{HidaSetting}$, surjectivity is the hypothesis on the
tested lattice map, its ends are crystalline by I82W2, and its
level-one unramified case, jumps and counts are the clauses above.
Both $H$ and $s_a+S_0$ are at most $h_B$:
for a thick letter use $s_a\le B+H$, for a nonthick one use
$s_a+S_0\le2b'\le b$. Thus $M_a+h_a\le 2h_B$. Together with
$b-A\le B\le b-A+1$ and the witness $E_a$ at each nonthick
letter, this proves $\mathrm{WordBox}$. The preceding numerical
clauses and I82W2 give all of $\mathrm{WordBase}$.
% END I82W3

% BEGIN I82W4
\paragraph{I82W4. The long-orbit word in row B.}
The frozen row B condition is $p<\ell_A$ and
$4L_w<\operatorname{ord}_H(p)$. Indeed arithmetic Frobenius
on the real $2^{e+2}$ cyclotomic subfield is represented by $p$
in the units modulo $\{\pm1\}$. Its order is the least $d>0$
with $p^d\equiv\pm1\bmod 2^{e+2}$, exactly the source's
$\operatorname{ord}_H(p)$. This projected order divides the
order of the second component of $\psi$, which divides $n_w$.
Let $j$ be the least integer with $\ell_A\le(p-1)^j$.
Since $11\le p<\ell_A$, $j=\lceil\log_{p-1}\ell_A\rceil
=j_A\ge2$. The parameter choice provides such a power with
$2j+2\le L_w$ at its chosen exponent, hence also at the least one.
In particular $j+1\le n_w$. If $C_2\ge0$, the box estimate gives
\[
 C_2(j+1)(W_0+S_0+H+1)\le C(j+1)b'\le b-b'\le B .
\]
The middle inequality uses $C(L_w+1)b'\le b$ and $C\ge100$.
For $C_2<0$ the left hand side is already nonpositive.

To see the cluster condition, write the additive second coordinate
of the letter at position $i$ as $y+i v$. A thick place has
coordinate $u_0+s a_0$, $s\in T_2$ in the host notation.
The generator $a_0$ has order $n_0$ and the host digits satisfy
$4L_w s<n_0$. If three positions $i_1<i_2<i_3$ in one interval
of length $L_w$ were thick, their three digits $s_1,s_2,s_3$
would be distinct since $\operatorname{ord}(v)>4L_w$.
Eliminating $v$ gives
\[
 n_0\mid(i_3-i_1)(s_2-s_1)-(i_2-i_1)(s_3-s_1).
\]
The absolute value is strictly below $n_0$ by the four digit
bounds. Equality to zero contradicts the host's no-relation
property for three distinct digits with distinct nonzero
coefficients $i_3-i_1,i_2-i_1$ bounded by $L_w$.
Thus every cyclic window of length $L_w$ has at most two marks.

Start a cluster at each marked letter with no mark in the preceding
$j$ positions. Take as its second letter the first mark in the next
$j$ positions, or use offset zero for a singleton if there is none.
If a mark has a predecessor within $j$, that predecessor must start
a cluster: another predecessor would give three marks within
$2j+1\le L_w$. The next mark is then the one we started from, by
the same three-mark test. Every mark is covered once, since two
distinct starts within $j$ contradict the definition of a start.
All intermediate letters of a pair and the $j$ positions before
its start are unmarked. The bound $2j<n_w$ permits cyclic indexing
without identifying these positions. In particular for the second letter
$t$ of a pair of offset $d$, the only mark among $t-k$ for
$1\le k\le d+j$ is at $k=d$: the first $d-1$ slots are
intermediate, and the last $j$ precede its start. This proves the native
$\mathrm{Clusters}(j_A)$ and hence, with I82W3,
$\mathrm{LocalRepresentationBoundCondition}$ at each row B point.
% END I82W4

% BEGIN I82W5
\paragraph{I82W5. The alternating words and the frozen family conclusion.}
The sign of $f_p$ on $\mathbb Q(\sqrt{-d})$ is the quadratic
symbol $(-d/p)=-1$. Its inverse $\psi$ has the same sign. The
orbit therefore has even length and changes orientation at every
step, including the wraparound. There is $d_0\in\{0,1\}$ for
which up means $a+d_0$ is odd.

If $\psi=\mathrm{cc}$, every letter has the same underlying real
place, so they are all thick or all nonthick. Suppose
$\psi\ne\mathrm{cc}$. If the second-component order exceeds
$4L_w$, the three-digit argument of I82W4 applies. Since $L_w\ge3$
by the parameter condition on word length, three consecutive thick
letters are impossible and at least one letter is nonthick. If
the second-component order is $m_v\le4L_w$, two thick hits with
digits $s,s'$ give $(i-i')v=(s-s')a_0$, hence
$n_0\mid m_v(s-s')$. The size bound on the digits forces $s=s'$.
The digit and the fixed translate specify the whole thick real
place. Consecutive letters with the same real place differ by
$1$ or $\mathrm{cc}$, both excluded here by the sign and the
case assumption. Thus this case even excludes two consecutive
thick letters. All cases with a nonthick letter satisfy
$\mathrm{WordsNT}$ by I82W3.

In the all-thick case put $M=\max s_a$, $Y=M+S_0+H$.
From I82W1--I82W3,
$B\ge100$, $2H+S_0\le B/10$ (use $H+S_0\le2b'$ and $B\ge99b'$),
and $B\le s_a,M\le B+H$. Therefore $100\le Y\le1.1B$.
For $p\ge11$, $(2p-1)/(p-1)^2\le21/100$ and
$\log_p(2Y)\le 2Y/100+4$, by $\log z\le z-1$ applied to
$z=2Y/100$, $\log 100<5$, and $\log p>1$. Thus
\[
 2a_R(p,Y)+H+1
 =2\{Y(2p-1)/(p-1)^2+4+\log_p(2Y)\}+H+1
 \le .506B+17+H\le B\le s_a .
\]
Likewise $(M+S_0)/2\le s_a$. This supplies both large-size
inequalities at the exact $h_{\rm par}=M+S_0$. The level-one
coefficient field contains a primitive $p$-th root by goodness.
I82W2 gives the required embedding-by-embedding small multiplicities
of every tested $T_2$. Hence this alternative gives
$\mathrm{WordBaseC}$ and $\mathrm{WordsAT}(p,d_0)$.

In particular the word in I82W2--I82W3 satisfies
$\mathrm{LocalRepresentationCongruenceCondition}$ on the entire
inert set, including the complement of row B where it is requested.
I82W4 supplies the other implication. Both use the original $X.r_i$,
$X.\rho_i$, arbitrary $\pi\in Z_{\rm good}(i)$, every $w\mid p$
and every good $L'$ at each twisted inert $i$. This is exactly
$\mathrm{JointParametersHubConditions}(L,HW,X,R)$ at the fixed
row constants on the late family. The local cohomological lifting
and length estimates that consume these word data are separate
obligations in the middle-range estimate.
% END I82W5
